% ============================================================
%  ch2_background  二、背景（原书 pp. 2230--2235）
% ============================================================

\section{背景}

\subsection{镜像对称}

\origpage{2230}对我们来说，镜像对称有三个相关的方面：拓扑镜像对称、同调镜像对称，以及 Strominger--Yau--Zaslow 镜像对称。

简而言之，镜像对称把一个 Calabi--Yau 流形 $X$（复结构 $I_X$，Kähler 形式 $\omega_X$，特别地也就是辛形式）的复几何，与同维的镜像 Calabi--Yau 流形 $(Y,I_Y,\omega_Y)$ 的辛几何联系起来。这类对应最早的 $3$ 维例子 \cite{greene-plesser} 出现在弦论与镜像对称中，其陈述是：以 $(X,\omega_X)$ 为目标的某个二维 $A$ 型非线性 $\sigma$ 模型的物理，与以 $(Y,I_Y)$ 为目标的二维 $B$ 型非线性 $\sigma$ 模型的物理相同。

对这一镜像对称的数学侧面的研究，一直是现代辛几何与复代数几何的中心课题之一。镜像关系的第一条数学侧面，是 Hodge 数相等：$h^{n-p,q}(X)=h^{p,q}(Y)$。我们称之为\textbf{拓扑镜像对称}。

1994 年，Kontsevich \cite{kontsevich} 提出了\textbf{同调镜像对称}的表述
\[
Fuk(X,\omega_X)\cong D^b(Y,I_Y),
\]
即 $X$ 的 Fukaya 范畴与 $Y$ 上凝聚层的导出范畴等价。这就为“镜像对称应当是什么意思”给出了一个深刻的数学猜想。

90 年代早期，\cite{candelas-etal1,batyrev1,batyrev-borisov} 中给出了若干猜想中的镜像对的构造。1996 年，Strominger--Yau--Zaslow \cite{strominger-etal} 提出一种办法，从 Calabi--Yau $3$ 维流形 $X$ 自身的几何出发构造它的镜像。他们论证，$X$ 与 $Y$ 应当纤维化在同一个底 $B\cong S^3$ 之上：
\begin{equation}\label{syz}
\begin{array}{ccccc}
X & & & & Y\\
 & \searrow & & \swarrow & \\
 & & B & &
\end{array}
\end{equation}
并且一般纤维是互为对偶的特殊拉格朗日 $3$-环面。这里，$L\subset X$ 是\emph{特殊拉格朗日}的，意思是 $L$ 是拉格朗日的（$\omega_X|_L=0$），并且 Calabi--Yau 体积形式的虚部也在 $L$ 上消失（$\Im(\Omega_X)|_L=0$）。反过来，$Y$ 应当看成 $Fuk(X,\omega_X)$ 中某些对象的模空间，这些对象一般由配上 $U(1)$-局部系统的特殊拉格朗日 $3$-环面给出。Strominger--Yau--Zaslow 的这一图景在数学上的实现由 Gross--Siebert \cite{gross-siebert1} 推进，已取得许多成果，新近更有突破 \cite{gross-siebert2,gross-etal}。

在互为镜像的 Calabi--Yau $3$ 维流形之间，对偶特殊拉格朗日纤维化 \eqref{syz} 的完整构造至今仍然缺失。Langlands 对偶群的 Higgs 丛模空间则是一个自然的例子：在它上面，Hitchin 系统自动给出这样的对偶特殊拉格朗日纤维化。不过，这与原始的 \cite{strominger-etal} 在几何情境上并不相同。

\subsection{几何 Langlands 对应}

在 \cite{drinfeld1}、\cite{laumon1}、\cite{beilinson-drinfeld} 的工作中，提出了 Langlands 对应的一个几何版本。回忆一下：数域上约化群 $\G$ 的 Langlands 纲领（入门性的介绍见 \cite{gelbart,frenkel}）把 $\G$ 的自守数据——例如 $\SL_2$ 的模形式——与某个 Langlands 对偶群 $\G^L$ 的谱数据——例如 $\SL_2$ 情形下椭圆曲线上同调里的一个 Galois 表示——联系起来。这些猜想也可以对另一种整体域来表述，即有限域上曲线的函数域；在后一种情形，因为可以有效地使用曲线的代数几何，猜想变得较为可行。这里我们要考虑的是这个纲领更加几何化的版本，即复数域上曲线的函数域的情形。

在 $\C$ 上，复约化群 $\G$ 的 Langlands 对偶 $\G^\vee:=\G^L$ 很容易构造。复约化群的分类由它们的根数据 $(X,\Phi,X^\vee,\Phi^\vee)$ 给出：一个秩为 $n$ 的格 $X$，以及对偶格 $X^\vee$ 中的根系 $\Phi\subset X$ 与余根系 $\Phi^\vee\subset X^\vee$，并满足若干条件；每个秩为 $n$ 的复约化群 $\G$ 都附有这样一份数据。$\G$ 的 Langlands 对偶就是以对偶根数据 $(X^\vee,\Phi^\vee,X,\Phi)$ 为根数据的约化群 $\G^\vee$。例如 $\GL_n^\vee\cong \GL_n$，$\SL_n^\vee\cong \PGL_n$。

对光滑射影曲线 $C$，Beilinson--Drinfeld \cite{beilinson-drinfeld} 的几何 Langlands 对应提出：由 $C$ 上的一个 $\G$-局部系统——Galois 表示的几何类比——出发，构造丛的模叠 $Bun_{\G^\vee}$ 上的一个 holonomic $D$-模——自守形式的几何类比。这个构造的主要性质是：该 holonomic $D$-模必须是某些 Hecke 算子的特征层。\cite{beilinson-drinfeld} 对 $C$ 上一类特定的 $\G$-局部系统，即所谓的 opers，完成了这一构造。

下面我们会看到，Beilinson--Drinfeld 的图景也可以理解为一种增强镜像对称，它发生在平坦 $\G$-联络的模空间 $\M_\DR$（它是超 Kähler 的，从而特别是 Calabi--Yau 的）与平坦 $\G^\vee$-联络的模空间 $\M^\vee_{\DR}$ 之间。

\subsection{Langlands 对偶 Hitchin 系统的 SYZ 镜像对称}

\cite{hausel-thaddeus} 的出发点是这样一个观察：当 $\G=\SL_n$、$\G^\vee=\PGL_n$ 时，$\G$-Higgs 丛的 Hitchin 系统 $h:\M_\Dol\to \A$ 与 $\G^\vee$-Higgs 丛的 Hitchin 系统 $\M^\vee_\Dol\to \calA^\vee$ 有同一个底 $\calA\cong \calA^\vee$，而在图
\begin{equation}\label{syzdol}
\begin{array}{rcccl}
\M_\Dol & & & & \M^\vee_\Dol\\
 & {}_{h} \searrow & & \swarrow_{h^\vee} & \\
 & & \calA & &
\end{array}
\end{equation}
中，一般纤维是互为对偶的阿贝尔簇。由于 Hitchin 系统完全可积，这些纤维关于全纯辛形式 $\omega_\C:=\omega_J+i\omega_K$ 是全纯拉格朗日的，即沿纤维有 $\omega_\C=0$，从而 $\omega_J=\omega_K=0$。

同一些映射在复结构 $J$ 下给出
\begin{equation}\label{syzdr}
\begin{array}{rcccl}
\M_\DR & & & & \M^\vee_\DR\\
 & {}_{h} \searrow & & \swarrow_{h^\vee} & \\
 & & \calA & &
\end{array}
\end{equation}
而此时一般纤维是特殊拉格朗日纤维化。这意味着 $\omega_J$ 在纤维上消失，因此纤维是拉格朗日的；此外 $\omega_K$ 以及 $J$-全纯 Calabi--Yau 形式 $(\omega_K+i \omega_I)^{2d}$ 的虚部也沿纤维消失。

为了得到镜像对称的一个可以在数学上检验的形式，\cite{hausel-thaddeus} 在互为镜像的 Calabi--Yau 流形 $\M_{\DR}$ 与 $\M^\vee_{\DR}$ 之间表述了\textbf{拓扑镜像对称}。

\subsection{Langlands 对偶 Hitchin 系统的拓扑镜像对称}

\origpage{2233}为表述 \cite{hausel-thaddeus} 的这个镜像对称版本，我们需要把模空间说得更精确些。对 $\SL_n$，取 $\M:=\M_\Dol$ 为 $C$ 上秩 $n$、$E$ 固定、行列式线丛固定且次数为 $1$、Higgs 场无迹（$\Phi\in H^0(C;\End_0(E)\otimes K)$）的稳定 Hitchin 对的模空间。对 $\PGL_n$，考虑 $\Gamma:=\Jac_C[n]$——即 $C$ 上 $n$ 阶线丛构成的群——在 $\SL_n$ 的模空间 $\M$ 上的作用，并令 $\M^\vee:=\M/\Gamma$。于是 $\M$ 是光滑的拟射影簇，而 $\M^\vee$ 是拟射影的轨形。在我们的讨论中，还需要模空间上的额外扭结构，其形式是一个 gerbe
\[
\alpha\in H^2(\M^\vee,\U(1))\cong H^2(\M,\U(1))^\Gamma,
\]
它可以用 $\M$ 上的万有丛构造出来。

于是可以对光滑的 $\M$ 定义某些混合 Hodge 数
\[
h^{p,q}(\M)=h^{p,q}(H^{p+q}_c(\M)),
\]
并对轨形 $\M^\vee=\M/\Gamma$ 定义 $\alpha$-扭的弦性 Hodge 数：
\[
h_{st,\alpha}^{p,q}(\M^\vee):=\sum_{\gamma\in \Gamma} h^{p-F(\gamma);q-F(\gamma)}\bigl(H^{p+q-2F(\gamma)}(\M^\gamma;L_{\alpha})^\Gamma\bigr),
\]
其中 $F(\gamma)$ 是费米子移位（fermionic shift），由 $\gamma$ 在 $\M$ 的某个 $\gamma$-不动点处的切空间上的作用定义。

有了这些，我们就可以表述拓扑镜像对称猜想：

\begin{conjecture}[\cite{hausel-thaddeus}]\label{tms}
Hodge 数相等：
\[
h^{p,q}(\M)=h_{st,\alpha}^{p,q}(\M^\vee).
\]
\end{conjecture}

要注意，按目前的表述，这个猜想是关于 Higgs 模空间 $\M=\M_{\Dol}$ 的 Hodge 数的。然而 \cite{hausel-thaddeus} 证明了 $\M_\Dol$ 与 $\M_{\DR}$ 的 Hodge 数相同，所以 \eqref{tms} 同时也是在说所提议的镜像 $\M_{\DR}$ 与 $\M^\vee_{\DR}$ 的（混合）Hodge 数相等。这样，我们就可以把猜想 \ref{tms} 解释为我们的 SYZ 镜像对 $\M_{\DR}$ 与 $\M^\vee_{\DR}$ 之间的拓扑镜像对称。

猜想 \ref{tms} 在 \cite{hausel-thaddeus} 中对 $\SL_2$ 与 $\SL_3$ 获证。新近，$\SL_n$ 的一般情形在 2020 年由 Gr\"ochenig--Wyss--Ziegler \cite{groechenig-etal1} 解决。他们用的是算术 $p$ 进积分技术；这一技术由 Denef--Loeser \cite{denef-loeser} 首创，Batyrev \cite{batyrev2} 曾用它检验通常镜像对称中的一些拓扑镜像对称猜想。

在 \cite{hausel-gths} 中，我们注意到一个耐人寻味的相似之处。Ng\^{o} 在证明 Langlands 纲领的基本引理 \cite{ngo} 时，把基本引理归结为：有限域上某些奇异 Hitchin 纤维的点数相等，进而归结为 Hitchin 纤维化某些奇异纤维的特定 Hodge 数相等。在 \cite{hausel-gths}[\S 5.4] 中我们论证，Ng\^{o} 的上同调结果是拓扑镜像对称猜想沿 Hitchin 纤维化的相对版本；我们还提出了一个用 Ng\^{o} 的技术推出拓扑镜像对称的策略。这一设想最近由 Maulik--Shen 在 \cite{maulik-shen} 中于 2020 年完成，他们用 Ng\^{o} 的技术给出了猜想 \ref{tms} 的一个新证明。

最后，2020 年 Gr\"ochenig--Wyss--Ziegler 在 \cite{groechenig-etal2} 中把他们从 \cite{groechenig-etal1} 得到的 $p$ 进积分技术推广到一般约化群 $\G$ 的 Higgs 模空间，并由此得到了 Ng\^{o} 上同调结果的新证明。

\subsection{几何 Langlands 作为增强同调镜像对称}\label{gls}

\origpage{2234}2007 年，Kapustin--Witten \cite{kapustin-witten} 提出了一整套细致的思想，其要旨是从物理上把几何 Langlands 对应推导为一种增强镜像对称。他们论证：某个四维 $N=4$ 超对称 Yang--Mills 理论中研究得很充分的 $S$ 对偶（又称电磁对偶，或 Montonen--Olive 对偶 \cite{montonen-olive}），约化到二维之后给出一种增强镜像对称，而后者又反过来蕴含了 \cite{beilinson-drinfeld} 所表述的几何 Langlands 对应。

在这一二维约化中，Montonen--Olive 对偶成为两个 $\sigma$ 模型之间的等价：一个是以复曲线 $C$ 上平坦 $\G$-联络的模空间 $\M_\DR$ 为目标的 $B$ 型 $\sigma$ 模型，另一个是以 $C$ 上平坦 $\G^\vee$-联络的模空间 $\M^\vee_\DR$ 为目标的 $A$ 型 $\sigma$ 模型。于是，两个理论中边界条件构成的范畴应当等价：
\begin{equation}\label{hmsgl}
\mathscr{S}:D^b(\M_\DR)\simeq Fuk(\M_\DR^\vee),
\end{equation}
这可以解释为把 Kontsevich 的同调镜像对称猜想应用于镜像对 $\M_\DR$ 与 $\M_\DR^\vee$。

Kapustin--Witten \cite{kapustin-witten} 指出：由于目标空间是超 Kähler 的，又由于这些范畴源于四维理论，这个范畴的等价比通常的同调镜像对称带有更多的结构与更多的对称性；后者只来自两个二维 $\sigma$ 模型之间的等价。他们用这些额外的想法，从 $\M_\DR$ 中的一个平坦联络 $\alpha$（用它的摩天大楼层 $\calO_\alpha\in D^b(\M_\DR)$ 来表示）出发，构造出 $Fuk(\M^\vee_\DR)$ 中的一个元素，并把它解释为我们曲线 $C$ 上 $\G^\vee$-丛的模空间上的一个 $D$-模。Hecke 特征层的性质则反过来由四维来源的额外对称性推出。

首先，由于目标空间是超 Kähler 的，Kapustin--Witten 讨论的是结构更丰富的膜（也就是边界条件）：他们提出，一个膜相对于三个复结构 $I$、$J$、$K$ 中的每一个，都应当是 $A$ 型或 $B$ 型。这样他们就考虑超 Kähler 流形上 $(B,A,A)$ 型、$(A,B,A)$ 型、$(A,A,B)$ 型与 $(B,B,B)$ 型的膜。例如，一个 $(B,A,A)$ 型膜可以是一个 $I$-全纯的、关于 $\omega_\C=\omega_J+i\omega_K$ 为拉格朗日的子簇，再配上其上的一个局部系统；而一个 $(B,B,B)$ 型膜则应当是一个超 Kähler 子流形，再配上它上面某个丛上的超全纯联络。

在其框架中，\cite{kapustin-witten} 论证：超 Kähler 流形 $\M$ 上一个 $(B,A,A)$ 型膜的镜像（即 $S$ 对偶）应当是 $\M^\vee$ 上的一个 $(B,B,B)$ 型膜。特别地，如果只看复结构 $J$（即 $\M_{\DR}$ 的那个），我们就得到：$A$ 型膜的镜像应当是 $B$ 型膜。

镜像关系还要更微妙一些 \cite[Table 2 p74]{kapustin-witten}：复结构 $J$ 下的 $A$ 模型应当镜像于复结构 $K$ 下的 $B$ 模型，而由 \eqref{JK} 后者又给出我们这里的说法。这个更精细的镜像对称版本，把复结构 $I$ 下的一个 $B$ 型膜，匹配到镜像上复结构 $I$ 下的另一个 $B$ 型膜。

后一种对应也被 Donagi--Pantev \cite{donagi-pantev} 表述为同调镜像对称 \eqref{hmsgl} 的经典极限，即 \eqref{hmsgl} 的一阶近似：
\begin{equation}\label{classlim}
\mathscr{S}:D^b(\M_\Dol)\simeq D^b(\M_\Dol^\vee).
\end{equation}
\cite{donagi-pantev} 首先验证：对每个约化群 $\G$，\eqref{syzdol} 的一般纤维确实是对偶阿贝尔簇。其次他们验证：一般来说，相对于 Hitchin 底的 Fourier--Mukai 变换 \cite{mukai} 给出一个形如 \eqref{classlim} 的等价，并且满足下面要讨论的 Hecke--Wilson 对称性的交换关系。可以预期，当把 \eqref{classlim} 延拓到模空间的某些奇点之上时，它必须作修改。

另一个研究方向——由 \cite{kapustin-witten} 对超 Kähler 膜的考察所激发——带来了对 $\M_\Dol$ 与 $\M_\DR$ 中拉格朗日子簇、以及 $\M$ 上超全纯层的新理解。例如 \cite{baraglia-schaposnik,garcia-prada-etal,hitchin-characteristic} 研究了全部四种类型中超 Kähler 膜的各种构造，并思考了它们的镜像应当是什么。特别地，Hitchin \cite{hitchin-characteristic} 提出了 $\M$ 上一个 $(B,A,A)$ 型膜与 $\M^\vee$ 上一个 $(B,B,B)$ 型膜组成的膜对，并证明它们在 Hitchin 底的一般点上是 Fourier--Mukai 对偶的。Hitchin 这一建议在 $\G=\GL_2$ 的情形是 \cite{hausel-mellit-pei} 的出发点；在那里，一个额外的结构，即 $\T$-作用，起了重要作用。这一观点下面还会更详细地解释。

其次，由于 \eqref{hmsgl} 的推导有四维来源，Kapustin--Witten 还考虑了这些范畴上来自四维理论中\emph{线算子}的额外对称性。具体地说，对 $\G$ 的表示 $\mu$，附有 Wilson 算子 $\mathscr{W}^\mu$；它通过张上万有 $\G$ 平坦联络所对应的 $\G$-丛在表示 $\mu$ 中的向量丛，作用在 $D^b(\M_\DR)$ 上。另一侧，对 $\G$ 的不可约表示 $\mu$，附有 Hecke 算子（物理学家称之为 't Hooft 算子）$\mathscr{H}^\mu$；它作用在 $\G^\vee$-丛的模空间上，进而作用在这些模空间上的 $D$-模上。\cite{kapustin-witten} 的额外对称性观察是：这些算子应当与 \eqref{hmsgl} 的镜像对称相交换。下面在 \S\ref{geosat} 中，我们会在 \cite{donagi-pantev} 的更对称的经典极限 \eqref{classlim} 里更详细地写出这些算子。

同调镜像对称 \eqref{hmsgl} 再加上这两项额外结构——镜像对称下超 Kähler 膜的匹配，以及 Wilson--Hecke 对称性——就是我们所称的\textbf{增强镜像对称}。它们超出了 Kontsevich 的通常同调镜像对称，源于 $S$ 对偶的超对称来源与四维来源。

